\section{Discussion}
\label{sec:discussion}

The recovery step asks whether the units a robust fit has discarded hide a group, by the
likelihood of Gaussian groups plus uniform noise, and it can follow any method that flags
units.
After TCLUST-REG at a generous level with reweighting it returned the groups of $15\%$ to $30\%$
that the level had trimmed away, and it reduced the cells with a mean accuracy below $0.85$, up
to $20\%$ of outliers, from $31$ of $89$ to four; the four that remain have three or four
groups, where the level discards part of several groups at once, and the step left them at
$0.83$ and $0.81$. After the trimmed cluster-weighted model it did the same.
No fixed level of TCLUST-REG between $0.25$ and $0.40$, with or without equal weights, was
the best choice on both sides of $20\%$ of outliers; the level $0.35$ came closest, and was
within $0.03$ of ESF in $83$ of the $89$ cells up to $20\%$ but lost $0.11$ to $0.13$ with four groups. ESF, the flagging procedure without a trimming level, is
the option when that knowledge is missing and the contamination is probably below a quarter.
On the three data sets ESF reached, without a level, the answers that TCLUST-REG gives at the
right level.

The step settles whether a poorly fitted set of units is contamination or a group, by the
likelihood of the model that the flagging rule already assumes, one candidate at a time and at
the given $K$; where the number of groups is in doubt, a scan over the number of
groups (Section~\ref{sec:recovery}) is the better tool. After a generous level the
step accepted narrow bands of uniform noise as a group at $30\%$ and $35\%$ of outliers, so it
should not follow such a level when a large fraction of uniform noise is possible.

The flagged fraction $\hat\alpha$ is the quantity a user of a level-free method reads first,
and it is a contamination estimate only under near-Gaussian errors. Proposition~\ref{thm:step}
says what it estimates then, the contamination plus about $1.2\%$ of clean units, and on the
simulated data with Gaussian errors it lay within about three percentage points of
$\varepsilon$ wherever ESF worked. On the tone data ESF flagged $0.25$ to $0.27$ of a benchmark
that is treated as clean, and on the taxi data $0.23$ against a contamination of $9.6\%$,
because the errors are far from Gaussian in both (Section~\ref{sec:realdata}); the fitted
lines were right in both cases. A flagged fraction well above any plausible contamination
therefore says that the errors are not Gaussian, and one well below it says that the fit has
broken down (Section~\ref{sec:constants}).

Set ESF against TCLUST-REG at level $0.30$ with reweighting and the recovery step. Up to
$20\%$ of outliers the two have the same mean accuracy over
the $89$ cells, $0.914$ against $0.915$, and the fixed level has no cell below $0.85$ where ESF
has two, at $0.78$ and $0.83$ (Table~\ref{tab:level}). ESF gives three things. It needs no
level, and the level $0.30$ is a choice that pays only because the contamination in those
cells is below it: beyond $25\%$ the same level falls below $0.85$ in $17$ of $34$ cells,
ESF in $13$. With four groups it is more accurate, $0.95$ against $0.88$
(Table~\ref{tab:mainlow}), because the trimming level discards part of several groups at once
and the estimated group weights make this worse; the level $0.25$ with equal weights removes
that loss ($0.96$) but not the loss beyond $25\%$. It costs three things: a subsample size $m$,
which encodes a bound on the smallest group proportion or, at $m=12$, gives up a little
accuracy (Section~\ref{sec:constants}); the cap of $n/3$ on the flagged set, which stops the
recovery step where the fit is in doubt and so limits the recovery of small groups to about
$20\%$ of outliers; and a dependence on the seed that the fit of TCLUST-REG with $1000$ starts does
not show, which the best-of-five rule reduces at five times the cost.

The choice among the methods of the study is as follows. If the contamination is known to be
below a fifth, TCLUST-REG at level $0.25$ with equal weights, reweighting and the recovery step
was the most accurate pipeline, with ESF within $0.03$ of it in $85$ of $89$ cells. If it may exceed a
quarter, use a level of $0.35$ or $0.40$ with reweighting, followed by the recovery step unless the
outliers may be spread uniformly. If nothing is known about it, ESF with the recovery step
needs no level and is accurate up to about a fifth, and a flagged fraction well below the
plausible contamination warns that it has broken down. With two groups and outliers in the response only, the Laplace
mixture was about as accurate as ESF up to $20\%$ of outliers and more robust beyond, but it
fails with three or four groups, with high-leverage points and with clustered outliers.

Five limits should be kept in mind. ESF breaks down when outliers are so frequent that clean
subsamples are rare in the first stages and a contaminated fit is taken for the reference; in
our designs this happened from $25\%$ of outliers with uniform background noise or
high-leverage points and from $30\%$ with three groups, while with two groups of equal size,
or with three covariates, it held up to $35\%$. Small groups are recovered only up to about
$20\%$ of outliers, and not on every data set at $20\%$ (Section~\ref{sec:gate2}). Outliers concentrated in a tight cloud
were fitted as part of a group from $20\%$ on; the recovery step repaired about two fifths of
these fits. If the error spread grows with the covariate, units with large covariate values
are flagged: with a spread growing fourfold over the design range the flagged fraction
exceeded the contamination by up to $6$ points (Plan~5), and in raw trade data the spread
grows much faster, so a transformation that stabilises it, such as the logarithm of
Section~\ref{sec:realdata}, should precede the fit. The candidate search of the recovery step, a line through $d$ random units among $500$
draws, is itself unstable when the group to be recovered is a small part of the flagged set:
on the fishery data, after TCLUST-REG at level $0.40$ with reweighting, the step found the
cheap group from three of ten starts (Section~\ref{sec:fishery}), so the step can miss a
group that it would accept. Its noise range is set by the most extreme units, so that
outliers far out make bands of noise easier to accept as groups (Section~\ref{sec:gate2}).
A result for the whole procedure would need to control which units are flagged, and we do
not have one (Section~\ref{sec:theory}).

The number of groups $K$ was taken as known. The Bayesian information criterion of the
Gaussian-lines-plus-noise model over several values of $K$ gives a way to choose it: in a
descriptive check (Supplement, Table~S18) it chose the true $K$ for at least $90\%$ of the
data sets with two or four groups and up to $20\%$ of outliers, but for only $50\%$ to $80\%$
with the three overlapping groups of D2, and it counted a tight cloud of outliers as a group;
on the fishery data it prefers $K=2$ to $K=1$ and $K=3$.
